\subsection{JSON Data Boundaries}

JSON is external text, not Python syntax. \texttt{json.loads}/\texttt{json.load} parse text into \texttt{dict}/\texttt{list}/scalars; \texttt{json.dumps}/\texttt{json.dump} serialize back.

\subsubsection{JSON as Text-Based Tree Serialization: Objects, Arrays, Strings, Numbers, Booleans, and Null}

JSON objects map to \texttt{dict}, arrays to \texttt{list}, \texttt{true}/\texttt{false} to \texttt{bool}, \texttt{null} to \texttt{None}. Until parsed, a JSON document is an ordinary \texttt{str}.

\begin{verbatim}
import json

json_text = '{"name": "Chuck", "level": 42, "active": true, "tags": ["fact"]}'
data = json.loads(json_text)
print(type(data), data["name"], data["tags"][0])
\end{verbatim}

\subsubsection{Loading JSON into Python Dictionaries and Lists with \texttt{json.load()} and \texttt{json.loads()}}

\texttt{json.loads} parses in-memory strings; \texttt{json.load} reads from a text file object. Both construct a new object graph on the heap.

%<<<PROGRAM:programs/jso01>>>
\subsubsection{Writing Python Structures Back to JSON with \texttt{json.dump()} and \texttt{json.dumps()}}

Serialization walks Python objects and emits UTF-8 text with JSON lexical rules (double-quoted keys/strings).

\subsubsection{JSON Type Mapping Boundaries: \texttt{None} vs. \texttt{null}, \texttt{dict} vs. Object, \texttt{list} vs. Array}

Only JSON-native types round-trip cleanly. Python \texttt{set}, \texttt{datetime}, and arbitrary class instances require custom encoders or pre-conversion.
